
\subsubsection{6.1 Root Cause Analysis}

The hypothesis failure has a clear mechanistic explanation: feature saturation in pretrained models.

CLIP ViT-B/16 was trained on approximately 400 million image-text pairs. Both background (land/water) and bird type are elementary visual concepts well within CLIP's training distribution. Consequently:

\begin{enumerate}
\item No emergence dynamics exist to measure
\item Linear probes converge immediately regardless of regularization strength
\item CV captures only sampling noise, not differential learning
\end{enumerate}

\subsubsection{6.2 Theoretical Interpretation}

The negative result clarifies a distinction between two phenomena:

\begin{itemize}
\item \textbf{Feature emergence}: A training-time phenomenon occurring as a model learns representations
\item \textbf{Feature separability}: A property of the final representation
\end{itemize}

C-sweep probing on frozen features measures separability, not emergence. The hypothesis conflated these by assuming that regularization strength (C) in convex optimization would serve as a proxy for training epochs in non-convex neural network optimization.

\subsubsection{6.3 Limitations}

\begin{enumerate}
\item \textbf{Single hypothesis tested}: Only the existence hypothesis (H-E1) was evaluated. The proposed intervention mechanism (gradient regularization based on CV) remains untested.
\end{enumerate}

\begin{enumerate}
\item \textbf{Single dataset and extractor}: Results are from Waterbirds with CLIP ViT-B/16. Other datasets (CelebA, ColoredMNIST) and feature extractors (DINOv2, ImageNet-pretrained ResNet) were not evaluated.
\end{enumerate}

\begin{enumerate}
\item \textbf{C-sweep as epoch proxy}: Using regularization strength as an epoch proxy may fundamentally fail to capture learning dynamics. Convex optimization does not exhibit the staged learning that neural network training does.
\end{enumerate}

\begin{enumerate}
\item \textbf{Two-sample AUC limitation}: With only two feature types (background and bird type), AUC is binary (0 or 1). Marginal CV differences combined with directional reversal yield the worst possible AUC.
\end{enumerate}

\subsubsection{6.4 Implications for Future Work}

The negative result identifies necessary conditions for emergence-based detection to succeed:

\begin{enumerate}
\item \textbf{Training-time measurement}: CV must be computed during model training, not on frozen pretrained representations
\end{enumerate}

\begin{enumerate}
\item \textbf{Unfrozen representations}: The model must be learning features during the measurement period
\end{enumerate}

\begin{enumerate}
\item \textbf{Alternative signals}: Loss curves, gradient norms, or other training-time metrics may provide emergence signals that post-hoc probing cannot
\end{enumerate}

